
%%%%%%%%%%%%%%%%%%%%%%%%%%%%%%%%%%%%%
%%      Fundamentação Teórica      %%
%%%%%%%%%%%%%%%%%%%%%%%%%%%%%%%%%%%%%

\chapter{Fundamentação Teórica}
    \label{cha:fundamentacao-teorica}

    \section{Análise e visualização de dados}

    A Análise de dados é uma tarefa complexa, mas de fundamental importância para a tomada de decisões \cite{guimaraes2004sistema}. Ela pode ser descrita como um processo de transformação de dados em informações que tenham valor. Com a crescente produção e coleta de dados, pelo surgimento do IOT e dos próprios usuários, esta tarefa tornou-se de suma importância para a evolução da Web como a conhecemos.

    O processo de análise de dados pode ser dividido em três etapas:

    \begin{itemize}
        \item \textbf{Pré-processamento}: como formatação, remoção de dados faltantes e valores atípicos (\textit{Outliers}), normalização, etc. Por exeplo a remoção de \textit{Outliers}, problemas durante a coleta ou armazenamento.
        \item \textbf{Processamento}: pode ser realizada de diversas formas, como, por exemplo, a aplicação de métodos estatísticos ou algoritmos de aprendizado de máquina. Por exemplo: regressão, classificação, agrupamento, predição, otimização, entre outros.
        \item \textbf{Pós-processamento}: consiste na apresentação dos resultados obtidos, que pode ser por simples tabelas, gráficos, modelos matemáticos ou até valores numéricos. Por exemplo: visualização de dados em gráficos de séries históricas.
    \end{itemize}

    Por fim, a visualização de dados é uma tarefa crucial para a realização de todo o processo, pois é através dela que o usuário realiza a avaliação dos dados, identifica o pré-processamento necessário, as técnicas de processamento que serão aplicadas e verifica a qualidade dos resultados obtidos.

    \section{Aprendizado de Máquina}

    O aprendizado de máquina consiste em conjunto de algoritmos visando resolver um problema de forma que algoritmos clássicos não conseguem ou não operam eficientemente \cite{monard2003conceitos}. Eles são capazes de gerar modelos matemáticos com base em dados, que podem ser utilizados para prever o comportamento dos dados, agrupar os dados ou classificá-los.

    Os problemas que estes algoritmos resolvem podem ser divididos em algumas categorias, como, por exemplo: regressão, classificação, agrupamento, predição, detecção de anomalias, ranqueamento, otimização, entre outros.

    \begin{itemize}
        \item \textbf{Regressão}: Envolve encontrar uma função que descreve a relação entre uma variável dependente e uma ou mais independentes. A função encontrada pode ser usada para prever o valor da variável dependente para novos valores da variável independente.
        \item \textbf{Classificação}: Consiste em encontrar uma função que descreve a relação entre uma variável dependente e uma ou mais variáveis independentes. A função encontrada pode ser utilizada para prever o valor da variável dependente para novos valores da variável independente.
        \item \textbf{Agrupamento}: Encontrar grupos de dados que compartilham características semelhantes. Logo, pode-se usar esses grupos para realizar uma análise mais detalhada de seus dados.
        \item \textbf{Predição}: Compreende a tarefa de prever um valor para uma variável dependente a partir de um conjunto de variáveis independentes. Por exemplo, prever o preço de uma ação a partir de uma série histórica de preços.
        \item \textbf{Detecção de anomalias}: Trata-se de encontrar dados com características diferentes dos demais. Esses dados podem ser removidos ou estudados separadamente.
        \item \textbf{Ranqueamento}: Esta técnica tem a finalidade de ordenar os dados de acordo com uma métrica de relevância. Por exemplo, ranquear os resultados de uma busca.
        \item \textbf{Otimização}: Consiste em encontrar uma função com objetivo de maximizar ou minimizar um valor. Por exemplo, encontrar o melhor caminho para um problema de roteamento.
    \end{itemize}

    \setlength{\leftskip}{3cm}

        Cada algoritmo citado, e muitos outros, já estão implementados em bibliotecas de código, algumas de código aberto, como, por exemplo, o \textit{scikit-learn}, e outras de código fechado, como, por exemplo, o MATLAB. Estas bibliotecas são utilizadas para estudo e criação de plataformas para a realização de análises de dados.

    \setlength{\leftskip}{0pt}

    \section{Engenharia de Software}

    O desenvolvimento de um software é complexo e envolve várias etapas, como o levantamento de requisitos, modelagem do software, desenvolvimento, testes, validação, implementação entre outras. Essas etapas buscam garantir a produção e qualidade do software conforme as necessidades levantadas previamente e validadas \cite{pressman2021engenharia}.

    A engenharia de software surgiu para facilitar e padronizar os projetos e o desenvolvimento de software, bem como para aumentar a eficiência e a qualidade do software produzido. Para isso, a engenharia de software utiliza diversas técnicas e ferramentas, como por exemplo: Padrões de Projeto, Testes, entre outras.

    Desenvolvendo algumas dessas técnicas: Padrões de Projeto, são soluções comuns para vários problemas de projeto, de forma com que possam ser utilizados em contextos diversos \cite{pressman2021engenharia}. E os Testes, cujo objetivo é garantir a qualidade do software, através da verificação de suas funcionalidades.

    \section{Linguagens de Programação}

    A linguagem escolhida para esse projeto foi o Python, presente entre as linguagens mais utilizadas no mundo, com uma vasta comunidade e um número crescente de bibliotecas. A linguagem Python é de alto nível, multiparadigma, possui tipagem dinâmica e forte \cite{python2022}. Além disso, utiliza de uma sintaxe simples e fácil de aprender para pessoas que já possuem conhecimento em outras linguagens de programação, assim como para iniciantes \cite{van1995python}.

    A linguagem executa sobre uma máquina virtual, chamada de Python Virtual Machine (\textit{PVM}), responsável por interpretar o código-fonte e executar as instruções \cite{pvm-2017}. O Python usa a compilação \textit{just-in-time} (JIT), que compila o código-fonte para \textit{bytecodes} e depois para código de máquina, em tempo de execução \cite{pvm-2017}.

    Em um software Python o código é compilado para \textit{bytecodes}, uma linguagem intermediária, que logo a seguir é transpilado para código de máquina pela \textit{PVM}, transpilar é o processo de converter o código de um estado para outro com uma série de otimizações e correções.

    O \textbf{CPython} é a implementação mais popular da linguagem Python, a implementação oficial da linguagem, e é a implementação que será utilizada nesse projeto. Desenvolvido em C, o \textbf{CPython} é multiplataforma, ou seja, pode ser executado em qualquer sistema operacional e arquitetura. O \textbf{CPython} funciona na seguinte maneira a seguir:

    \begin{enumerate}
        \item O código-fonte é compilado para \textit{bytecodes}.
        \item O \textit{bytecodes} é interpretado pela \textit{PVM}.
        \item A \textit{PVM} executa o código.
        \item O código é compilado para código de máquina.
        \item O código de máquina é executado.
    \end{enumerate}

    O Python utiliza um sistema de gerenciamento de memória automático, responsável por gerenciar a alocação e desalocação de memória. O Python possui um coletor de lixo, responsável por liberar a memória que não está sendo utilizada \cite{pvm-2017}. A linguagem trata bem questões relacionadas a paralelismo, concorrência e processamento assíncrono, por bibliotecas como o \textit{asyncio}, a qual é uma biblioteca de código aberto, que implementa o modelo de programação assíncrona.

    Em relação à desempenho, a linguagem é relativamente lenta, quando comparado a opções como C e C++, mas possui uma comunidade ativa, com uma vasta gama de bibliotecas e \textit{frameworks}, esse é o principal \textit{trade-off} da linguagem. Logo assim, devido a sua fama, comunidade, rica biblioteca e a alta curva de aprendizado, o Python foi escolhido para o desenvolvimento desse projeto.
%provavelmente vão questionar alguns dos motivos.
    \section{Bibliotecas}

    Bibliotecas são conjuntos de códigos que podem ser importados para outros projetos, realizando tarefas específicas com resolução já implementadas. A seguir serão apresentadas as principais bibliotecas utilizadas nesse projeto:

    \subsection{Numpy}

    O Numpy é uma biblioteca fundamental para a computação científica. Ela possui implementações de vetores n-dimensionais em C, compiladas e executadas de maneira nativa, resultando em alto desempenho \cite{numpy}. Além disso, o Numpy já dispõe de uma vasta gama de funções matemáticas e operações de álgebra linear. Usando o Numpy é possível realizar, por exemplo, a normalização de um conjunto de dados com apenas uma linha de código.

    \subsection{Pandas}

    O Pandas, que no que lhe concerne recorre ao Numpy, é uma biblioteca \textit{open-source}, ou seja, é um projeto disponibilizado com código aberto, que estrutura os dados em tabelas (\textit{Dataframes}), que são estruturas multidimensionais de dados, utilizados para análise de dados \cite{pandas}. Por utilizar o Numpy, além de uma vasta gama de funções matemáticas, o Pandas também possui desempenho de alto nível. Com o Pandas é possível realizar várias operações com os dados, como ordenação, filtragem, entre outros.

    \subsection{Scikit-learn}

    O Scikit-learn, como o próprio nome sugere, é uma biblioteca de \textit{machine learning}, que possui implementações de algoritmos de aprendizado de máquina, estatísticos e de mineração de dados \cite{scikit-learn2022}. O \textit{Scikit-learn} possui um conjunto considerável de algoritmos de classificação, regressão, agrupamento, redução de dimensionalidade, entre outros. Além disso, ele possui um grupo de métricas de avaliação, utilizadas para avaliar o desempenho dos modelos criados. Utilizando a biblioteca é possível realizar, por exemplo, a criação de um modelo de árvore de decisão, um exemplo de algoritmo de aprendizado supervisionado, utilizado para classificação e regressão.

    \subsection{PyQt}

    O PyQt é a biblioteca responsável por criar as interfaces gráficas de usuário (GUI) do projeto. Sobe licença GPL, o PyQt é uma biblioteca de código aberto, que opera por \textit{wrapper's}, funções que encapsulam as bibliotecas pré-compiladas do Qt \cite{pyqt}. O Qt no que lhe concerne, é uma biblioteca, também de código aberto, que dispões de uma série de \textit{widgets}, layouts, entre outros componentes que podem ser facilmente utilizados para a criação de interfaces gráficas. Outra vantagem do PyQt é sua compatibilidade com vários sistemas operacionais, como Linux, Windows, Mac OS, entre outros.

    \subsection{Seaborn/Matplotlib}

    O Seaborn é a biblioteca gráfica escolhida nesse projeto para tratar a visualização dos dados. Ela utiliza o Matplotlib como base, fornecendo várias interfaces de alto nível para a criação e personalização de gráficos \cite{seaborn}. O Seaborn possui um conjunto de possibilidades como, por exemplo, a criação de gráficos de dispersão, gráficos de linha, gráficos de barras, entre outros. Utilizando ela é possível embutir os gráficos gerados na interface gráfica criada com o PyQt.

    % ÚLTIMO TÓPICO
    \section{Trabalhos Relacionados}
    \label{sec:trabalhos-relacionados}

    [PENDENTE]

